\documentclass{book}

\usepackage{../themes/dark_theme}

\begin{document}
	
	\pagenumbering{Roman}
    \begin{titlepage}
        %\tikz[remember picture,overlay] \node[opacity=0.5,inner sep=0pt] at (current page.center){\includegraphics[width=\paperwidth,height=\paperheight]{"arriere_plan.png"}};

        \vspace*{9em}{\centering\Huge Algèbre\par}
        \vspace{2em}
        {\hfill\itshape Enzo S. Sanfilippo}

        \clearpage
	    \vspace*{\fill}
    \end{titlepage}

    \newpage
    \pagenumbering{roman}

    {
		\hypersetup{linkcolor=magenta}
		\tableofcontents
	}

    \newpage
	\pagenumbering{arabic}

	\chapter{Algèbre}{
		\newpage

		\section{Théorèmes de \texorpdfstring{\textsc{Sylow}}{Sylow}}{
			\noindent Soit $p \in \mathbb{P}$. \\
			\noindent Soit $n \in \mathbb{N}$. \\
			\noindent Soit $k \in \mathbb{N}^{\star}$ tel que $p \wedge k = 1$. \\
			\noindent Soit $G$ un groupe fini d'ordre $p^{n} \times k$.

			\mathblock{theorem}{\textsc{Sylow} $1$}{
				Alors $G$ possède un sous-groupe $S$ tel que $|S| = p^{n}$.
			}

			Un tel sous-groupe est appelé $p$-sous-groupe de \textsc{Sylow}.

			\mathblock{definition}{Ensemble des $p$-sous-groupes de \textsc{Sylow}}{
				\begin{mathlist}
					\item On notera
						$$
							\mathrm{Syl}_{p}(G)
						$$
						l'ensemble des $p$-sous-groupes de \textsc{Sylow} de $G$.
					\item On notera
						$$
							n_{p} := \mathrm{Card} \ \mathrm{Syl}_{p}(G)
						$$
				\end{mathlist}
			}

			\mathblock{theorem}{\textsc{Sylow} $2$}{
				Soit $P \leqslant G$ un $p$-sous-groupe de $G$ ($p \ | \ |P|$). \\
				Soit $S \in \mathrm{Syl}_{p}(G)$ un $p$-sous-groupe de \textsc{Sylow}. \bigskip \\
				Alors pour tout $g \in G$,
				$$
					P^{g} \subset S
				$$
			}

			\mathblock{corollary}{Lien entres les $p$-sous-groupes de \textsc{Sylow}}{
				Soient $S, T \in \mathrm{Syl}_{p}(G)$. \bigskip \\
				Alors il existe $g \in G$ tel que
				$$
					S^{g} = T
				$$
			}

			\mathblock{theorem}{\textsc{Sylow} $3$}{
				On a
				$$
					n_{p} \ | \ k
				$$
			}

			\mathblock{theorem}{\textsc{Sylow} $4$}{
				On a
				$$
					n_{p} \equiv 1[p]
				$$
			}
		}
		\section{Produit de groupes}{
			\noindent Soit $G$ un groupe quelconque.

			\mathblock{lemma}{Critère pour multiplier des groupes}{
				Soient $A, B \leqslant G$ deux sous-groupes distingués. \bigskip \\
				Alors
				$$
					AB \leqslant G \iff AB = BA
				$$
			}

			On rappelle que
			$$
				AB := \{ a \times b, a \in A, b \in B \}
			$$

			\mathblock{corollary}{Critère plus faible}{
				Soit $A \triangleleft G$. \\
				Soit $B \leqslant G$. \bigskip \\
				Alors
				$$
					AB \leqslant G
				$$
			}

			\mathblock{theorem}{\textsc{Frottini}}{
				Soit $N \triangleleft G$ fini. \\
				Soient $p \in \mathbb{P}, k \in \mathbb{N}^{\star}$ et $n \in \mathbb{N}^{\star}$ tels que $|N| = p^{n} k$ avec $p \wedge k = 1$. \\
				Soit $S \in \mathrm{Syl}_{p}(N)$. \bigskip \\
				Alors
				$$
					G = N N_{G}(S)
				$$
			}
		}

		\section{Série de sous-groupes}{
			\noindent Soit $G$ un groupe fini.

			\mathblock{definition}{Chaîne de sous-groupes}{
				Soit $n \in \mathbb{N}^{\star}$. \\
				Soit $(H_{i})_{i \in I}$ une famille de sous-groupes de $G$. \bigskip \\
				On dit que $(H_{i})_{0 \leq i \leq n}$ est une \textit{chaîne de sous-groupes} lorsque
				$$
					\forall i \in [\![0, n-1]\!], H_{i} \leqslant H_{i + 1}
				$$
			}

			\mathblock{definition}{Série de sous-groupes}{
				Soit $n \in \mathbb{N}^{\star}$. \\
				Soit $(H_{i})_{i \in I}$ une famille de sous-groupes de $G$. \bigskip \\
				On dit que $(H_{i})_{0 \leq i \leq n}$ est une \textit{série de sous-groupes} lorsque
				\begin{mathlist}
					\item $(H_{i})_{0 \leq i \leq n}$ est une chaîne de sous-groupes.
					\item $H_{0} = \{ e \}$.
					\item $H_{n} = G$.
				\end{mathlist}
			}

			\mathblock{definition}{Série normale / sous-normale}{
				Soit $n \in \mathbb{N}^{\star}$. \\
				Soit $(H_{i})_{i \in I}$ une famille de sous-groupes de $G$. \bigskip \\
				On dit que
				\begin{mathlist}
					\item $(H_{i})_{0 \leq i \leq n}$ est une \textit{série normale} lorsque
						$$
							\forall i \in [\![0, n]\!], H_{i} \triangleleft G
						$$
					\item $(H_{i})_{0 \leq i \leq n}$ est une \textit{série sous-normale} lorsque
						$$
							\forall i \in [\![0, n-1]\!], H_{i} \triangleleft H_{i+1}
						$$
				\end{mathlist}
			}
		}
	}
\end{document}
